\paragraph{D4m (Proved: a modular Igusa degeneration with flat
$\operatorname{Pic}^{\tau}$ but incompatible Picard varieties).}
This example is independent of the full versal-ring calculation.
Let $k$ be algebraically closed of characteristic two, let $R=W(k)$,
and take both $E$ and $F$ to be the elliptic curve
\[
 Y^2Z+XYZ=X^3+XZ^2
\]
over $R$, with origin $O=[0:1:0]$ and $P=[0:0:1]$.
The discriminant is $-63$, a unit of $R$, so these are smooth elliptic
schemes. The point $P$ has order two and specializes to a nonzero
point; hence the special fibres are ordinary. Use subscripts to
distinguish $P_E$ and $P_F$.

Let $w\in\mathbb Z_2\subset R$ be the unique root congruent to zero
modulo two of
\[
 2w^2-w+8=0.
\]
It exists uniquely by Hensel's lemma, because the derivative is odd;
also $w\in8R$. The point
\[
 Q_F=[-2:1:w]\in F(R)
\]
lies on $F$, is fixed by elliptic negation
$[X:Y:Z]\mapsto[X:-Y-X:Z]$, and is nonzero on the generic fibre.
Thus it has order two generically and as an $R$-section, whereas
$(Q_F)_k=O$. It is the nontrivial generic section of the multiplicative
two-torsion subgroup. No choice of a root of a nonunit in $R$ is
being assumed here; the displayed Hensel equation defines $Q_F$.

Define two commuting involutions on $E\times_R F$ by
\[
 h(x,y)=(x+P_E,y+Q_F),\qquad
 \tau(x,y)=(-x,y+P_F).
\]
They generate the constant group $\Gamma=(\mathbb Z/2)^2$.
The action is free. For $h$ this follows from its nonzero $E$-translation;
for $\tau$ and $h\tau$ it follows from their $F$-translations
$P_F$ and $P_F+Q_F$, both of which specialize to a nonzero point.
The zero locus of either section is closed in $\operatorname{Spec}R$
and misses its closed point, so it is empty. Consequently
\[
 \mathcal X=(E\times_R F)/\Gamma
\]
is smooth projective over $R$, and its quotient map is finite etale.
For quotient existence one may use Stacks, Lemma~66.14.2, tag07S7;
an invariant ample tensor product descends to show projectivity.
Its special fibre is the ordinary Igusa surface
\[
 \mathcal X_k\simeq
 \bigl((E_k/\langle P_E\rangle)\times F_k\bigr)
 /\langle(-1,+P_F)\rangle.                              \tag{M1}
\]

\emph{The torsion Picard scheme is flat.}
First divide by $h$. Since $h$ is translation by the etale subgroup
$\langle(P_E,Q_F)\rangle$, the quotient
$A=(E\times F)/\langle h\rangle$ is an abelian scheme.
The copy of $F$ embeds in $A$, because that subgroup meets
$0\times F$ trivially. Its quotient is $\overline E=E/\langle P_E\rangle$.
Duality gives
\[
 0\longrightarrow\overline E^{\vee}\longrightarrow A^{\vee}
 \longrightarrow F^{\vee}\longrightarrow0.              \tag{M2}
\]
The linear part $\sigma$ of the remaining deck involution acts by
$-1$ on the first term and by $+1$ on the last term. If
$G=(A^{\vee})^{\sigma^{\vee}}$, the same elementary lifting argument
works over every base change: lift a point of $F^{\vee}$ to $A^{\vee}$,
then subtract half of its failure of invariance in
$\overline E^{\vee}$, taking the half fppf locally. Therefore
\[
 0\longrightarrow\overline E^{\vee}[2]\longrightarrow G
 \longrightarrow F^{\vee}\longrightarrow0               \tag{M3}
\]
is fppf exact. The map $G\to F^{\vee}$ is a finite locally free
torsor of degree four.

For the remaining double cover $A\to\mathcal X$, cyclic descent
gives
\[
 0\longrightarrow\mu_2\longrightarrow
 \operatorname{Pic}^{\tau}_{\mathcal X/R}
 \longrightarrow G\longrightarrow0.                    \tag{M4}
\]
Indeed an invariant algebraically trivial line bundle can be
linearized fppf locally: its chosen isomorphism with the deck
pullback has square a base unit, and adjoining a square root of
that unit corrects the cocycle. The ambiguity is exactly a
character of the cyclic group, namely $\mu_2$. The descended
class lies in $\operatorname{Pic}^{\tau}$ because its square is
the norm of an algebraically trivial bundle, hence is in the
neutral Picard component on each fibre. Conversely every torsion
Neron--Severi class pulls back into $A^{\vee}$, since the
Neron--Severi group of an abelian variety is torsion free.
These are fppf-sheaf statements on arbitrary $R$-schemes.
They also prove representability of the torsion Picard functor
here by finite-flat torsor descent from $G$.
Combining (M3) and (M4),
$\operatorname{Pic}^{\tau}_{\mathcal X/R}\to F^{\vee}$
is finite locally free of degree eight. In particular the full
torsion Picard scheme is flat over $R$.

\emph{The Picard varieties cannot specialize as an abelian subscheme.}
Let $f:F\to\mathcal X$ send $y$ to the class of $(0,y)$ and use
$f(0)$ as base point. For either geometric fibre $s$, every pointed
map $\mathcal X_s\to Z$ to an abelian variety pulls back to a
homomorphism $u_E+u_F:E_s\times F_s\to Z$. Invariance under $\tau$
forces $2u_E=0$ and $u_F(P_F)=0$. Since homomorphisms of abelian
varieties have no additive torsion, $u_E=0$. Invariance under $h$
then forces $u_F(Q_F)=0$. Conversely these conditions suffice for
descent. The Albanese variety and the map induced by $f_s$ are thus
\[
 \operatorname{Alb}(\mathcal X_s)=
 F_s/\langle P_F,Q_F\rangle.
\]
On the generic fibre the two distinct nonzero rational two-torsion
points $P_F,Q_F$ are independent, so the subgroup is all of $F[2]$
and the map has degree four. On the special fibre $Q_F=0$, so its
degree is two. Dualizing, restriction along $f$ on the fibrewise
Picard varieties has degrees
\[
 \deg f_\eta^*=4,\qquad \deg f_k^*=2.                    \tag{M5}
\]
If a closed abelian subscheme of
$\operatorname{Pic}^{\tau}_{\mathcal X/R}$ had these Picard varieties
as all its fibres, restriction along $f$ would be an isogeny of
abelian schemes with the two degrees in (M5). An isogeny of abelian
schemes is finite locally free, so its degree is locally constant.
This contradiction proves the asserted specialization failure
inside an actual modular Igusa degeneration, even though
$\operatorname{Pic}^{\tau}$ itself is flat.

The finite calculation records the discriminant, the projective
point and two-torsion identities, and Hensel approximations. See
\texttt{check\_igusa\_trait.py} and its saved output. The calculation
does not infer flatness from tangent dimensions and does not use
the full-hull assertion.
